
\documentclass{article}
\usepackage{colortbl}
\usepackage{makecell}
\usepackage{multirow}
\usepackage{supertabular}

\begin{document}

\newcounter{utterance}

\centering \large Interaction Transcript for game `cladder', experiment `full\_v1.5\_default', episode 2287 with selfrefine{-}llama{-}3.1{-}8b{-}instruct.
\vspace{24pt}

{ \footnotesize  \setcounter{utterance}{1}
\setlength{\tabcolsep}{0pt}
\begin{supertabular}{c@{$\;$}|p{.15\linewidth}@{}p{.15\linewidth}p{.15\linewidth}p{.15\linewidth}p{.15\linewidth}p{.15\linewidth}}
    \# & \multicolumn{2}{c}{Player} && \multicolumn{2}{c}{Game Master} \\
    \hline

    \theutterance \stepcounter{utterance}  
    & & & \multicolumn{4}{p{0.6\linewidth}}{
        \cellcolor[rgb]{0.9,0.9,0.9}{
            \makecell[{{p{\linewidth}}}]{
                \texttt{\tiny{[P1$\langle$GM]}}
                \texttt{You are an expert in causal inference. The following question is not a typical commonsense query, but rather a meticulously designed question created by a professor specializing in causal inference, intended to assess the students' mastery of the course content.} \\
\texttt{Imagine a self{-}contained, hypothetical world with only the following conditions, and without any unmentioned factors or causal relationships: Smoking has a direct effect on effort and college admission. Effort has a direct effect on college admission. The overall probability of smoking is 93\%. The probability of nonsmoking and college admission is 3\%. The probability of smoking and college admission is 62\%. Is the chance of college admission smaller when observing smoking?} \\
\texttt{Start your answer with 'Yes' or 'No', followed by additional reasoning or evidence to support your explanation.} \\
            }
        }
    }
     \\ \\

    \theutterance \stepcounter{utterance}  
    & & \multicolumn{4}{p{0.6\linewidth}}{
        \cellcolor[rgb]{0.95,0.95,0.95}{
            \makecell[{{p{\linewidth}}}]{
                \texttt{\tiny{[P1$\leftrightarrow$P]}}
                \texttt{{[}DRAFT{]}} \\
\texttt{No.} \\
\\ 
\texttt{To answer this question, we need to apply the concept of causal inference, specifically the backdoor criterion or the frontdoor adjustment. However, in this case, we can use a simpler approach based on the given probabilities.} \\
\\ 
\texttt{Let's denote the event of smoking as S, the event of nonsmoking as NS, and the event of college admission as A. We are given the following probabilities:} \\
\\ 
\texttt{{-} P(S) = 0.93} \\
\texttt{{-} P(NS) = 1 {-} P(S) = 0.07} \\
\texttt{{-} P(A|NS) = 0.03} \\
\texttt{{-} P(A|S) = 0.62} \\
\\ 
\texttt{We want to determine if the chance of college admission is smaller when observing smoking. This is equivalent to comparing P(A|S) and P(A|NS).} \\
\\ 
\texttt{To do this, we can use the law of total probability to calculate P(A) and compare it with P(A|S) and P(A|NS).} \\
\\ 
\texttt{P(A) = P(A|S) * P(S) + P(A|NS) * P(NS)} \\
\texttt{= 0.62 * 0.93 + 0.03 * 0.07} \\
\texttt{= 0.5771 + 0.0021} \\
\texttt{= 0.5792} \\
\\ 
\texttt{Now, we can compare P(A) with P(A|S) and P(A|NS). We see that P(A|S) = 0.62 is greater than P(A) = 0.5792, and P(A|NS) = 0.03 is less than P(A) = 0.5792.} \\
\\ 
\texttt{However, the question asks if the chance of college admission is smaller when observing smoking. This is equivalent to asking if P(A|S) < P(A|NS). Since P(A|S) = 0.62 > P(A|NS) = 0.03, the answer is No.} \\
\\ 
\texttt{In other words, the chance of college admission is not smaller when observing smoking. In fact, the probability of college admission is higher when observing smoking compared to nonsmoking.} \\
            }
        }
    }
    & \\ \\

    \theutterance \stepcounter{utterance}  
    & & \multicolumn{4}{p{0.6\linewidth}}{
        \cellcolor[rgb]{0.95,0.95,0.95}{
            \makecell[{{p{\linewidth}}}]{
                \texttt{\tiny{[P1$\leftrightarrow$P]}}
                \texttt{{[}SELF{-}FEEDBACK{]}} \\
\texttt{FORMAT: ok} \\
\texttt{RULES: ok} \\
\texttt{STRATEGY: ok} \\
\texttt{REASON: The answer correctly applies the law of total probability to compare the probabilities of college admission under different conditions, and provides a clear explanation for the conclusion.} \\
            }
        }
    }
    & \\ \\

    \theutterance \stepcounter{utterance}  
    & \multicolumn{4}{p{0.6\linewidth}}{
        \cellcolor[rgb]{0.561,0.737,0.561}{
            \makecell[{{p{\linewidth}}}]{
                \texttt{\tiny{[P1$\rangle$GM]}}
                \texttt{No.} \\
\\ 
\texttt{The answer correctly applies the law of total probability to compare the probabilities of college admission under different conditions, and provides a clear explanation for the conclusion. } \\
\\ 
\texttt{P(A) = P(A|S) * P(S) + P(A|NS) * P(NS)} \\
\texttt{= 0.62 * 0.93 + 0.03 * 0.07} \\
\texttt{= 0.5771 + 0.0021} \\
\texttt{= 0.5792} \\
\\ 
\texttt{Now, we can compare P(A) with P(A|S) and P(A|NS). We see that P(A|S) = 0.62 is greater than P(A) = 0.5792, and P(A|NS) = 0.03 is less than P(A) = 0.5792.} \\
\\ 
\texttt{However, the question asks if the chance of college admission is smaller when observing smoking. This is equivalent to asking if P(A|S) < P(A|NS). Since P(A|S) = 0.62 > P(A|NS) = 0.03, the answer is No.} \\
\\ 
\texttt{In other words, the chance of college admission is not smaller when observing smoking. In fact, the probability of college admission is higher when observing smoking compared to nonsmoking.} \\
            }
        }
    }
    & & \\ \\

    \theutterance \stepcounter{utterance}  
    & & & \multicolumn{2}{p{0.3\linewidth}}{
        \cellcolor[rgb]{0.9,0.9,0.9}{
            \makecell[{{p{\linewidth}}}]{
                \texttt{\tiny{[GM$|$GM]}}
                \texttt{No} \\
            }
        }
    }
    & & \\ \\

    \theutterance \stepcounter{utterance}  
    & & & \multicolumn{2}{p{0.3\linewidth}}{
        \cellcolor[rgb]{0.9,0.9,0.9}{
            \makecell[{{p{\linewidth}}}]{
                \texttt{\tiny{[GM$|$GM]}}
                \texttt{no} \\
            }
        }
    }
    & & \\ \\

    \theutterance \stepcounter{utterance}  
    & & & \multicolumn{2}{p{0.3\linewidth}}{
        \cellcolor[rgb]{0.9,0.9,0.9}{
            \makecell[{{p{\linewidth}}}]{
                \texttt{\tiny{[GM$|$GM]}}
                \texttt{game\_result = WIN} \\
            }
        }
    }
    & & \\ \\

\end{supertabular}
}

\end{document}
